\documentclass[11pt]{amsart}
\usepackage[utf8]{inputenc}
\usepackage{xcolor}
\usepackage{amsmath,amssymb,amsthm}
\usepackage{booktabs}
\usepackage[expansion=false]{microtype}
\setlength{\emergencystretch}{2em}
\usepackage[margin=1in]{geometry}
\usepackage[colorlinks=true,linkcolor=blue!50!black,citecolor=blue!50!black,urlcolor=blue!50!black]{hyperref}

\newtheorem{theorem}{Theorem}[section]
\newtheorem{lemma}[theorem]{Lemma}
\newtheorem{proposition}[theorem]{Proposition}
\newtheorem{corollary}[theorem]{Corollary}
\theoremstyle{definition}
\newtheorem{definition}[theorem]{Definition}
\theoremstyle{remark}
\newtheorem{remark}[theorem]{Remark}

\newcommand{\F}{\mathbb F}
\newcommand{\HD}[1]{D^{[#1]}}
\newcommand{\rank}{\operatorname{rank}}
\newcommand{\ord}{\operatorname{ord}}
\newcommand{\Span}{\operatorname{span}}
\newcommand{\dd}{\,\mathrm d}
\renewcommand{\Re}{\operatorname{Re}}

\title[Hasse jets and an elliptic spectral bound]{Hasse jets, finite-field point counts, and an elliptic spectral bound:\\
a reproducible Stepanov argument}
\author{Vincent P.\,P. No\"el}
\address{Strasbourg, France}
%\email{} % to be completed by the author
\date{October 2026. Expository preprint, not peer reviewed.}
\subjclass[2020]{Primary 11G20; Secondary 11G05, 11Y16, 14G15}
\keywords{Stepanov's method, Hasse derivatives, finite fields, elliptic curves, point counts, Riemann--Roch, exact certificates}

\begin{document}

\begin{abstract}
We present a classical auxiliary-polynomial argument for cubic curves over finite fields of odd characteristic, together with exact reproducible certificates. Hasse derivatives give a jet-rank bound over every $\F_q$, $q=p^n$, including orders at least $p$. A polynomial space with distinct leading degrees, followed by multiplicity counting and quadratic twisting, yields the uniform bound $|\#C(\F_q)-(q+1)|\le10\sqrt q$, with asymptotic constant $2\sqrt6$. For a fixed smooth cubic over $\F_p$, rationality of its zeta function is established separately by Riemann--Roch. The bound on all extensions then forces both inverse roots to have modulus $\sqrt p$, recovering Hasse's sharp constant $2$ and the elliptic Riemann hypothesis. Five exact polynomial certificates and twelve direct extension-field counts accompany the proof. A singular cubic shows why the elliptic rationality hypothesis is essential. This is an expository reconstruction of Stepanov's method and classical elliptic-curve results; no new theorem or transfer to the Riemann zeta function is claimed.
\end{abstract}

\maketitle

\section{Introduction and scope}\label{sec:intro}

A point-counting estimate with a non-optimal constant can still determine the exact size of the inverse roots of a zeta function. The necessary ingredients are a bound uniform over \emph{all} extensions of the base field, a rational expression for the zeta function, and a known product of its inverse roots. In genus one, the product is $p$ and a root of modulus larger than $\sqrt p$ would dominate its partner exponentially.

This note makes that chain explicit for $y^2=f(x)$ with $f$ cubic and the characteristic odd. Its counting part uses the classical method of Stepanov~\cite{Ste69,Bom74}. Its rationality input is the standard Riemann--Roch computation for elliptic curves~\cite[I, \S4; IV, \S9]{Mil20}. Neither Hasse's bound nor a purity theorem is imported into the argument. These results and methods are classical; their presentation here is not a claim of historical novelty.

There are two points worth keeping separate in an implementation. First, the auxiliary space is a space of \emph{polynomials}, not of polynomial functions modulo $x^q-x$: a nonzero kernel vector must give a genuinely nonzero polynomial. Second, the relation between ordinary and Hasse derivatives uses $k!$ only when $k<p$. We instead compute Hasse derivatives directly, so the proof works for $M<q$ even when $M\ge p$. The extension-field certificates test precisely this distinction.

The argument supplies a concrete example of a structural rank bound, as opposed to a numerical positivity test. It does not construct an analogous space or operator for the zeta function of the integers. Section~\ref{sec:jets} proves the rank bound, Section~\ref{sec:count} the uniform point count, Sections~\ref{sec:rationality}--\ref{sec:spectral} the elliptic spectral consequence, and Section~\ref{sec:certificates} records the exact computations and their reproduction.

\section{Hasse jets and the auxiliary polynomial space}\label{sec:jets}

Throughout the counting argument, $q=p^n\ge5$ is an odd prime power, $f\in\F_q[x]$ has degree exactly three, and
\[
e=\frac{q-1}{2},\qquad g=f^e,\qquad
S=\{u\in\F_q:f(u)\ne0,\ g(u)=1\}.
\]
No squarefreeness assumption is needed until Section~\ref{sec:rationality}.

\begin{definition}
For a polynomial $A$, the Hasse derivative $\HD{k}A$ is defined by
\[
A(x+h)=\sum_{k\ge0}\HD{k}A(x)h^k.
\]
The jet of order $M-1$ at $u$ is $(\HD{0}A(u),\ldots,\HD{M-1}A(u))$.
\end{definition}
Directly from the definition,
\[
\HD{k}x^a=\binom ak x^{a-k}\quad(k\le a),\qquad
\HD{k}(AB)=\sum_{\ell=0}^k\HD{k-\ell}A\,\HD{\ell}B;
\]
the first expression is zero for $k>a$. All binomial coefficients are reduced in characteristic $p$. The product rule has no extra factor $\binom k\ell$. The first $M$ jet coordinates at $u$ vanish exactly when $(x-u)^M$ divides $A$, by expanding $A(u+h)$.

\begin{lemma}\label{lem:Pk}
For every $k\ge0$ there is a polynomial $P_k$ such that
\[
f^k\HD{k}g=gP_k,\qquad \deg P_k\le2k.
\]
One may take
\[
P_k=\sum_{\ell=0}^{\min(e,k)}
 \binom e\ell f^{k-\ell}
 [h^k]\bigl(f(x+h)-f(x)\bigr)^\ell.
\]
The degree of the zero polynomial is understood as $-\infty$.
\end{lemma}
\begin{proof}
Expand $f(x+h)^e$ binomially. Its coefficient of $h^k$ uses only terms $\ell\le k$, since $f(x+h)-f(x)$ is divisible by $h$. Multiplication by $f^k$ gives the displayed identity after factoring out $f^e$. The expression for $P_k$ is polynomial because $k-\ell\ge0$.

Write $f(x+h)-f(x)=c_1(x)h+c_2(x)h^2+c_3h^3$, where $\deg c_i\le3-i$. A term contributing to $[h^k](f(x+h)-f(x))^\ell$ has degree at most $3\ell-k$ in $x$. After multiplication by $f^{k-\ell}$, its degree is at most $3(k-\ell)+3\ell-k=2k$. This proof never divides by $k!$.
\end{proof}

Let
\[
J\ge1,\qquad 1\le K\le\frac{q-3}{2},\qquad 1\le M<q,
\]
and work in $\F_q[x]$, without quotienting by $x^q-x$. Set
\[
W=\Span_{\F_q}\{x^{jq+a},\,x^{jq+a}g:
                  0\le j<J,\ 0\le a<K\}.
\]

\begin{theorem}[Auxiliary-polynomial rank bound]\label{thm:rank}
The space $W$ has dimension $2JK$. The jet map
\[
\mathcal J_M:W\longrightarrow\F_q^{M|S|},\qquad
A\longmapsto(\HD{k}A(u))_{\substack{u\in S\\0\le k<M}}
\]
satisfies
\begin{equation}\label{eq:rank}
\rank\mathcal J_M
 \le\sum_{k=0}^{M-1}\min\{|S|,K+J-1+2k\}
 \le R:=M(K+J-1)+M(M-1).
\end{equation}
If $2JK>R$, there is a nonzero $A\in W$ vanishing to order at least $M$ at every point of $S$, and
\begin{equation}\label{eq:SU}
|S|\le U:=\left\lfloor\frac DM\right\rfloor,\qquad
D=(J-1)q+K-1+\frac{3(q-1)}2.
\end{equation}
\end{theorem}
\begin{proof}
\emph{Dimension.} The leading degrees in the two families are $jq+a$ and
\[
jq+a+\frac{3(q-1)}2=(j+1)q+\frac{q-3}{2}+a.
\]
Their residues modulo $q$ lie in disjoint intervals
\[
[0,K-1],\qquad
\left[\frac{q-3}{2},\frac{q-3}{2}+K-1\right]\subset[0,q-1].
\]
All $2JK$ leading degrees are distinct, so the polynomials are independent. Their maximum degree is $D$.

\emph{Frobenius and the rank of each block.} Since
\[
(x+h)^{jq}=(x^q+h^q)^j,
\]
we have $\HD{i}x^{jq}=0$ for $0<i<q$. Fix $k<M$ and multiply the jet row at $u\in S$ by the nonzero scalar $f(u)^k$. Using $u^{jq}=u^j$, the first family's rows are evaluations of
\[
x^jf^k\HD{k}x^a.
\]
For the second family, Hasse's product rule, Lemma~\ref{lem:Pk} and $g(u)=1$ give evaluations of
\[
x^j\sum_{\ell=0}^k f^{k-\ell}\HD{k-\ell}x^a\,P_\ell.
\]
Each nonzero term in either expression has degree at most $j+a+2k\le K+J-2+2k$. Thus the block of order $k$ factors through evaluation of polynomials of that degree and has rank at most $\min(|S|,K+J-1+2k)$. Stacking the blocks proves~\eqref{eq:rank}. This sum is an upper bound; it does not assert that the blocks are independent.

\emph{Existence and multiplicities.} If $\dim W>R$, the kernel of $\mathcal J_M$ has dimension at least $2JK-R>0$. Independence of the chosen polynomials ensures that a nonzero vector in this kernel gives a nonzero polynomial $A$. Its zero jets imply
\[
\prod_{u\in S}(x-u)^M\mid A.
\]
The factors are coprime for distinct $u$. Therefore $M|S|\le\deg A\le D$, proving~\eqref{eq:SU}.
\end{proof}

The bound $R$ and the parameters guaranteeing existence do not depend on a prior count of $S$. An exact witness computed using $S$ is an additional certificate for a particular instance.

\section{Point counts, twisting, and a uniform constant}\label{sec:count}

An explicit family of parameters will suffice. For $r\ge1$ with $6r^2-2r+5\le q$, choose
\[
J=r+1,\qquad M=2r,\qquad K=3r^2-r+1.
\]
The condition on $r$ ensures $K\le(q-3)/2$ and $M<q$. Direct calculation gives $2JK-R=2$ and
\begin{equation}\label{eq:Ufamily}
U=\left\lfloor
 \frac q2+\frac{3(q-1)}{4r}+\frac{3r-1}{2}
 \right\rfloor.
\end{equation}

Let $z$ be the number of distinct roots of $f$ in $\F_q$, so $0\le z\le3$. The projective model of $y^2=f(x)$ has one point at infinity and
\[
N_q=2|S|+z+1.
\]
For a non-residue $\nu\in\F_q^\times$, the twisted cubic $\nu f$ has the same roots and
\[
|S_{\nu f}|=q-z-|S|.
\]
It has the same bound $U$, since the degree remains three. Applying the bound to both cubics gives
\begin{equation}\label{eq:twist}
2q-z-2U+1\le N_q\le2U+z+1,\qquad
|N_q-(q+1)|\le2U-q+z.
\end{equation}
These are counts of the projective model, even for a singular cubic. We apply the elliptic zeta function only to smooth cubics below.

\begin{proposition}\label{prop:uniform}
For every odd prime power $q$, including $q=3$, and every cubic $f$,
\[
|N_q-(q+1)|\le10\sqrt q.
\]
The displayed family also gives the asymptotic bound
\[
|N_q-(q+1)|\le2\sqrt6\sqrt q+O(1),
\]
with an absolute implicit constant.
\end{proposition}
\begin{proof}
For $3\le q<29$, the quadratic character, extended by $\chi_q(0)=0$, gives
\[
N_q-(q+1)=\sum_{u\in\F_q}\chi_q(f(u)),
\]
whose absolute value is at most $q<10\sqrt q$. This treats the small fields without requiring any admissible auxiliary parameters.

For $q\ge29$ put $r=\lfloor\sqrt{(q-5)/6}\rfloor\ge2$. This is admissible. If $v=\sqrt{(q-5)/6}\ge2$, then $r\ge v/2$ and $r\le\sqrt{q/6}$. Equations~\eqref{eq:Ufamily}--\eqref{eq:twist}, with $z\le3$, yield
\[
|N_q-(q+1)|\le\frac{3(q-1)}{2r}+3r+2.
\]
Since $q/(q-5)\le29/24$,
\[
|N_q-(q+1)|
\le\left(\frac32\sqrt{29}+\sqrt{\frac32}+
                 \frac2{\sqrt{29}}\right)\sqrt q
<10\sqrt q.
\]
The parenthetical constant is less than $9.674$.

Finally $r=\sqrt{q/6}+O(1)$ in~\eqref{eq:Ufamily}, so $U=q/2+\sqrt6\sqrt q+O(1)$. Equation~\eqref{eq:twist} proves the asymptotic statement. Both constants are independent of $p$, the extension degree, and the coefficients of $f$.
\end{proof}

The constant $10$ is convenient rather than optimal. The following argument needs uniformity, not an improvement of this intermediate constant.

\section{Rationality obtained independently by Riemann--Roch}\label{sec:rationality}

Fix an odd prime $p$ and a squarefree cubic $f\in\F_p[x]$. Its smooth projective model $E$, with the point $O$ at infinity, is an elliptic curve. Write $N_n=\#E(\F_{p^n})$ and $t=p+1-N_1$.

We use Riemann--Roch in genus one: a divisor $D$ of positive degree $d$ has $\dim L(D)=d$~\cite[I, Thm.~4.13, pp.~35--36]{Mil20}. This input has no hypothesis on point-counting errors or zeta-root moduli. For completeness, the rationality computation and its descent to $\F_p$ are included.

\begin{proposition}[Classical elliptic rationality]\label{prop:rationality}
One has
\[
Z(E,T)=\frac{1-tT+pT^2}{(1-T)(1-pT)}
      =\exp\left(\sum_{n\ge1}\frac{N_n}{n}T^n\right).
\]
If $(1-\alpha T)(1-\beta T)=1-tT+pT^2$, then
\begin{equation}\label{eq:spectrum}
\alpha+\beta=t,\qquad\alpha\beta=p,\qquad
N_n=p^n+1-\alpha^n-\beta^n.
\end{equation}
Also $Z(E,1/(pT))=Z(E,T)$. None of these statements assumes Hasse's bound.
\end{proposition}
\begin{proof}
Let $\operatorname{Pic}(E)$ denote divisors defined over $\F_p$ modulo divisors of nonzero functions in $\F_p(E)$. The map $P\mapsto[(P)-(O)]$ is a bijection $E(\F_p)\to\operatorname{Pic}^0(E)$. Indeed, for a divisor $D$ of degree zero, Riemann--Roch applied to $D+(O)$ provides a nonzero $\F_p$-rational section. It makes $D+(O)$ equivalent to an effective rational divisor of degree one, hence to $(P)$ for some $P\in E(\F_p)$. Uniqueness follows from $\dim L((P))=1$: a function with divisor $(Q)-(P)$ is constant, forcing $P=Q$. Translation by $d(O)$ identifies degree-zero classes with degree-$d$ classes.

There are therefore $N_1$ classes of each degree $d$. A class represented by $D$ of degree $d\ge1$ has $(p^d-1)/(p-1)$ effective divisors, obtained from the nonzero elements of $L(D)$ modulo $\F_p^\times$. If $a_d$ counts all effective divisors of degree $d$, then
\[
a_0=1,\qquad a_d=N_1\frac{p^d-1}{p-1}\quad(d\ge1).
\]
The Euler product on closed points is their generating function, so
\[
\begin{aligned}
Z(E,T)
 &=\sum_{d\ge0}a_dT^d\\
 &=1+\frac{N_1}{p-1}
       \left(\frac{pT}{1-pT}-\frac{T}{1-T}\right)
 =\frac{1-tT+pT^2}{(1-T)(1-pT)}.
\end{aligned}
\]
This is the standard proof of~\cite[IV, Prop.~9.9 and Thm.~9.10,
pp.~155--156]{Mil20}; it does not use the earlier Hasse theorem in that book.

If $b_d$ counts the closed points of degree $d$, then $N_n=\sum_{d\mid n}d b_d$. Taking the formal logarithm of $\prod_P(1-T^{\deg P})^{-1}$ gives $\sum_{n\ge1}N_nT^n/n$. Comparing it with the logarithm of the rational factorisation proves~\eqref{eq:spectrum}. No analytic convergence or purity is needed.

Finally, both the numerator and denominator transform by the same factor $1/(pT^2)$ under $T\mapsto1/(pT)$, proving the functional equation.
\end{proof}

At this stage $\alpha\beta=p$ gives the reciprocal pairing $\alpha\leftrightarrow p/\alpha$. It does not yet determine $|\alpha|$ or $|\beta|$.

\section{Uniform counting determines the elliptic spectrum}\label{sec:spectral}

\begin{theorem}[Classical elliptic spectral bound]\label{thm:spectral}
For the smooth cubic $E$ of Section~\ref{sec:rationality},
\[
|\alpha|=|\beta|=\sqrt p,\qquad
|N_n-(p^n+1)|\le2p^{n/2}\quad(n\ge1).
\]
The zeros of $\zeta_E(s):=Z(E,p^{-s})$ have real part $1/2$.
\end{theorem}
\begin{proof}
Squarefreeness persists under extension of the ground field. Proposition~\ref{prop:uniform}, applied to each $\F_{p^n}$, and~\eqref{eq:spectrum} give
\[
|\alpha^n+\beta^n|\le10p^{n/2}\quad\text{for every }n\ge1.
\]
Suppose, after interchanging the roots if necessary, that $A=|\alpha|>\sqrt p$. Then $|\beta|=p/A<\sqrt p$. The reverse triangle inequality implies
\[
10\ge p^{-n/2}|\alpha^n+\beta^n|
 \ge\left(\frac A{\sqrt p}\right)^n
       \left(1-\left(\frac p{A^2}\right)^n\right).
\]
The last expression tends to infinity, a contradiction. Thus neither modulus exceeds $\sqrt p$, and their product forces equality. The triangle inequality now gives the sharp constant $2$, including Hasse's bound at $n=1$.

In the variable $s$,
\[
\zeta_E(s)=
 \frac{(1-\alpha p^{-s})(1-\beta p^{-s})}
      {(1-p^{-s})(1-p^{1-s})}.
\]
A zero satisfies $p^s=\alpha$ or $\beta$, hence $p^{\Re s}=\sqrt p$. The denominator has poles only on the lines $\Re s=0,1$, so it does not cancel these zeros. Zeros and poles are periodic in the imaginary direction, with period $2\pi/\log p$.
\end{proof}

The passage from constant $10$ to constant $2$ is a consequence of rationality and the uniform estimate on the entire tower. Twelve finite checks alone would not supply this uniform estimate. The proof does.

\section{Exact certificates and direct counts}\label{sec:certificates}

The accompanying corpus release V6.35 contains five exact auxiliary-polynomial certificates and twelve direct point counts, all for $f=x^3+x+1$ except the final certificate, which uses the twist $2f$ over $\F_{125}$. The characteristic-$5$ and characteristic-$3$ examples include $M>p$.

\begin{table}[ht]
\centering\footnotesize
\setlength{\tabcolsep}{4pt}
\begin{tabular}{rcrrrrrrr}
\toprule
$q$ & $(J,K,M)$ & $|S|$ & $\dim W$ & $M|S|$ & $R$ &
 $\rank\mathcal J_M$ & $\deg A$ & $\min\ord_u A$\\
\midrule
27 & $(3,11,4)$ & 13 & 66 & 52 & 64 & 51 & 93 & 4\\
25 & $(7,9,6)$ & 13 & 126 & 78 & 120 & 78 & 136 & 6\\
49 & $(3,11,4)$ & 27 & 66 & 108 & 64 & 56 & 173 & 4\\
125 & $(5,45,8)$ & 52 & 450 & 416 & 448 & 400 & 701 & 8\\
$125^\ast$ & $(5,45,8)$ & 70 & 450 & 560 & 448 & 400 & 701 & 8\\
\bottomrule
\end{tabular}
\caption{Exact polynomial certificates. The starred row uses $2f$. Orders are taken at $u\in S$. $R$ is the bound independent of $S$; the sharper bound in~\eqref{eq:rank} is also checked.}
\label{tab:certificates}
\end{table}

Each stored certificate includes the irreducible field modulus, the coefficients of $f$, the points of $S$, a kernel vector, and the coefficients of the resulting nonzero polynomial $A$. Its multiplicities are checked by successive synthetic divisions by $x-u$, independently of the jet matrix used to construct it.

The $\F_{27}$ instance tests the Hasse calculation beyond the characteristic: $M=4>3$. Its $52$ formal conditions are already fewer than the $66$ dimensions, so existence in this particular example does not by itself demonstrate the advantage of the structural rank bound. The $\F_{125}$ twisted instance does: it has $560$ conditions on $450$ independent coefficients, yet $R=448<450$ guarantees existence. The computed rank is $400$ and the kernel has dimension $50$. Its witness has degree $701\le730$ and multiplicity at least $8$ at each of the $70$ points.

\begin{table}[ht]
\centering\small
\begin{tabular}{rrrrrr}
\toprule
$p$ & $n$ & $q=p^n$ & $N_n$ (direct) & $N_n$ (recurrence) &
 $N_n-(q+1)$\\
\midrule
3 & 1 & 3 & 4 & 4 & 0\\
3 & 2 & 9 & 16 & 16 & 6\\
3 & 3 & 27 & 28 & 28 & 0\\
3 & 4 & 81 & 64 & 64 & $-18$\\
5 & 1 & 5 & 9 & 9 & 3\\
5 & 2 & 25 & 27 & 27 & 1\\
5 & 3 & 125 & 108 & 108 & $-18$\\
7 & 1 & 7 & 5 & 5 & $-3$\\
7 & 2 & 49 & 55 & 55 & 5\\
7 & 3 & 343 & 380 & 380 & 36\\
11 & 1 & 11 & 14 & 14 & 2\\
11 & 2 & 121 & 140 & 140 & 18\\
\bottomrule
\end{tabular}
\caption{Twelve direct counts for $y^2=x^3+x+1$. Direct field enumeration is completed before comparison with the recurrence.}
\label{tab:counts}
\end{table}

The comparison recurrence is $T_0=2$, $T_1=t$ and $T_n=tT_{n-1}-pT_{n-2}$, giving $N_n=p^n+1-T_n$ from Proposition~\ref{prop:rationality}. The initial $t$ is obtained from the direct count over $\F_p$. This recurrence does not generate any of the direct counts in Table~\ref{tab:counts}.

\subsection*{Reproduction and implementation limits}
In the accompanying archive \texttt{programme-riemann}, work from \path{3_Preuves_et_calculs}. With the dependencies in the corpus requirements file installed, run
\begin{verbatim}
python src/v635_tour_et_spectre.py
python tests/run_all_tests.py
\end{verbatim}
The core implementation is \path{src/finite_field_jets.py}. The output files are \path{data/v635_certificats.json}, \path{data/v635_tour.json} and \path{data/v635_tour.csv}; the declared cases are in \path{data/v635_controles_avant_execution.json}. The proof is also recorded in \path{preuves/LEMME_TOUR_ET_SPECTRE_V635.md}.

Field elements are base-$p$ coefficient vectors in $\F_p[u]/(m(u))$, with a checked monic irreducible modulus. Their integer labels are encodings, not residues modulo the integer $p^n$. Polynomial coefficients are stored in increasing degree order. Matrix columns are ordered by $j$, then $a$, with the plain and mixed columns alternating. The computations use exact field arithmetic; no floating-point eigenvalue threshold enters a certificate. Runtime measurements and normalised ratios are ancillary, and need not reproduce byte for byte.

The present table-based field implementation is deliberately restricted to $q\le1024$. Its polynomial witnesses and direct counts are finite checks, not the proof of an assertion for all $q$. The ten dedicated V6.35 tests include Hasse orders beyond the characteristic, field construction, exact jet and multiplicity checks, twisting, and the separation between direct enumeration and the spectral recurrence.

\section{Why the smoothness hypothesis matters}\label{sec:limits}

A singular cubic supplies a useful counterexample to a careless spectral conclusion. Let $a$ be a non-square in $\F_p$ and take $f=x^2(x+a)$. For the singular projective model,
\[
N_n-(p^n+1)
 =\sum_{x\ne0}\chi_{p^n}(x+a)
 =-\chi_{p^n}(a)=-(-1)^n.
\]
Thus
\[
Z_{\mathrm{sing}}(T)=
 \frac{1+T}{(1-T)(1-pT)}.
\]
This count has even bounded error, but its zero $T=-1$ has modulus $1$. The numerator has degree one and does not have the elliptic product $\alpha\beta=p$. One must distinguish the singular model from its normalisation. The counting lemma is valid for both smooth and singular cubics; Theorem~\ref{thm:spectral} is elliptic and requires squarefreeness.

The argument does not cover characteristic $2$, and the displayed parameter family is not claimed to optimise all auxiliary spaces. Extending or optimising an implementation does not change the historical status of the underlying method.

For the zeta function over the integers, the missing step remains structural: an independent space, an operator producing a rank deficit, and a global control of zeros that can replace polynomial degree. These three objects have not been supplied by this note. Its conclusion is the classical elliptic Riemann hypothesis over finite fields.

\subsection*{Acknowledgements and disclosure}
This expository reconstruction, the extension-field implementation, and the manuscript were prepared with assistance from OpenAI's ChatGPT in the author's research corpus. Two internal constructive and critical readings were used; neither is an independent human referee report. The computations and mathematical sources are identified above. Final verification and responsibility for any submission rest with the author.

\begin{thebibliography}{99}
\bibitem{Ste69}
S. A. Stepanov, \emph{On the number of points of a hyperelliptic curve over a finite prime field},
Izv. Akad. Nauk SSSR Ser. Mat. 33 (1969), no.~5, 1103--1114.
\url{https://mathnet.ru/eng/im2197}.

\bibitem{Bom74}
E. Bombieri, \emph{Counting points on curves over finite fields (d'apr\`es S. A. Stepanov)},
S\'eminaire Bourbaki, exp.~430 (June 1973), in \emph{S\'eminaire Bourbaki 1972/73},
Lecture Notes in Mathematics 383, Springer, 1974, 234--241.
\url{https://www.numdam.org/item/SB_1972-1973__15__234_0/}.

\bibitem{Mil20}
J. S. Milne, \emph{Elliptic Curves}, second edition, World Scientific, 2021.
See Chapter I, Theorem~4.13 and formula~(8), pp.~35--36, and Chapter IV,
Proposition~9.9, Theorem~9.10 and Remark~9.11(a), pp.~155--156.
Author's text: \url{https://www.jmilne.org/math/Books/EC2.pdf};
mathematics site: \url{https://www.jmilne.org/math/}.
\end{thebibliography}

\end{document}

